%!TeX root=../princesstop.tex
\chapter{The Diamond Mines Again}

\begin{figure}[t!]
\centering
\includegraphics[width=\linewidth]{filigree}
\end{figure}

\lettrine[lines=5]{W}{hen} Sara entered the holly-hung schoolroom in the afternoon, she did so as the head of a sort of procession. Miss Minchin, in her grandest silk dress, led her by the hand. A manservant followed, carrying the box containing the Last Doll, a housemaid carried a second box, and Becky brought up the rear, carrying a third and wearing a clean apron and a new cap. Sara would have much preferred to enter in the usual way, but Miss Minchin had sent for her, and, after an interview in her private sitting room, had expressed her wishes.

<This is not an ordinary occasion,> she said. <I do not desire that it should be treated as one.>

So Sara was led grandly in and felt shy when, on her entry, the big girls stared at her and touched each other's elbows, and the little ones began to squirm joyously in their seats.

<Silence, young ladies!> said Miss Minchin, at the murmur which arose. <James, place the box on the table and remove the lid. Emma, put yours upon a chair. Becky!> suddenly and severely.

Becky had quite forgotten herself in her excitement, and was grinning at Lottie, who was wriggling with rapturous expectation. She almost dropped her box, the disapproving voice so startled her, and her frightened, bobbing curtsy of apology was so funny that Lavinia and Jessie tittered.

<It is not your place to look at the young ladies,> said Miss Minchin. <You forget yourself. Put your box down.>

Becky obeyed with alarmed haste and hastily backed toward the door.

<You may leave us,> Miss Minchin announced to the servants with a wave of her hand.

Becky stepped aside respectfully to allow the superior servants to pass out first. She could not help casting a longing glance at the box on the table. Something made of blue satin was peeping from between the folds of tissue paper.

<If you please, Miss Minchin,> said Sara, suddenly, <mayn't Becky stay?>

It was a bold thing to do. Miss Minchin was betrayed into something like a slight jump. Then she put her eyeglass up, and gazed at her show pupil disturbedly.

<Becky!> she exclaimed. <My dearest Sara!>

Sara advanced a step toward her.

<I want her because I know she will like to see the presents,> she explained. <She is a little girl, too, you know.>

Miss Minchin was scandalized. She glanced from one figure to the other.

<My dear Sara,> she said, <Becky is the scullery maid. Scullery maids—er—are not little girls.>

It really had not occurred to her to think of them in that light. Scullery maids were machines who carried coal scuttles and made fires.

<But Becky is,> said Sara. <And I know she would enjoy herself. Please let her stay—because it is my birthday.>

Miss Minchin replied with much dignity:

<As you ask it as a birthday favour—she may stay. Rebecca, thank Miss Sara for her great kindness.>

Becky had been backing into the corner, twisting the hem of her apron in delighted suspense. She came forward, bobbing curtsies, but between Sara's eyes and her own there passed a gleam of friendly understanding, while her words tumbled over each other.

<Oh, if you please, miss! I'm that grateful, miss! I did want to see the doll, miss, that I did. Thank you, miss. And thank you, ma'am,>—turning and making an alarmed bob to Miss Minchin—<for letting me take the liberty.>

Miss Minchin waved her hand again—this time it was in the direction of the corner near the door.

<Go and stand there,> she commanded. <Not too near the young ladies.>

Becky went to her place, grinning. She did not care where she was sent, so that she might have the luck of being inside the room, instead of being downstairs in the scullery, while these delights were going on. She did not even mind when Miss Minchin cleared her throat ominously and spoke again.

<Now, young ladies, I have a few words to say to you,> she announced.

<She's going to make a speech,> whispered one of the girls. <I wish it was over.>

Sara felt rather uncomfortable. As this was her party, it was probable that the speech was about her. It is not agreeable to stand in a schoolroom and have a speech made about you.

<You are aware, young ladies,> the speech began—for it was a speech—<that dear Sara is eleven years old today.>

<\textsc{dear} Sara!> murmured Lavinia.

<Several of you here have also been eleven years old, but Sara's birthdays are rather different from other little girls' birthdays. When she is older she will be heiress to a large fortune, which it will be her duty to spend in a meritorious manner.>

<The diamond mines,> giggled Jessie, in a whisper.

Sara did not hear her; but as she stood with her green-gray eyes fixed steadily on Miss Minchin, she felt herself growing rather hot. When Miss Minchin talked about money, she felt somehow that she always hated her—and, of course, it was disrespectful to hate grown-up people.

<When her dear papa, Captain Crewe, brought her from India and gave her into my care,> the speech proceeded, <he said to me, in a jesting way, <I am afraid she will be very rich, Miss Minchin.> My reply was, <Her education at my seminary, Captain Crewe, shall be such as will adorn the largest fortune.> Sara has become my most accomplished pupil. Her French and her dancing are a credit to the seminary. Her manners—which have caused you to call her Princess Sara—are perfect. Her amiability she exhibits by giving you this afternoon's party. I hope you appreciate her generosity. I wish you to express your appreciation of it by saying aloud all together, <Thank you, Sara!>>

The entire schoolroom rose to its feet as it had done the morning Sara remembered so well.

<Thank you, Sara!> it said, and it must be confessed that Lottie jumped up and down. Sara looked rather shy for a moment. She made a curtsy—and it was a very nice one.

<Thank you,> she said, <for coming to my party.>

<Very pretty, indeed, Sara,> approved Miss Minchin. <That is what a real princess does when the populace applauds her. Lavinia>—scathingly—<the sound you just made was extremely like a snort. If you are jealous of your fellow-pupil, I beg you will express your feelings in some more lady-like manner. Now I will leave you to enjoy yourselves.>

The instant she had swept out of the room the spell her presence always had upon them was broken. The door had scarcely closed before every seat was empty. The little girls jumped or tumbled out of theirs; the older ones wasted no time in deserting theirs. There was a rush toward the boxes. Sara had bent over one of them with a delighted face.

<These are books, I know,> she said.

The little children broke into a rueful murmur, and Ermengarde looked aghast.

<Does your papa send you books for a birthday present?> she exclaimed. <Why, he's as bad as mine. Don't open them, Sara.>

<I like them,> Sara laughed, but she turned to the biggest box. When she took out the Last Doll it was so magnificent that the children uttered delighted groans of joy, and actually drew back to gaze at it in breathless rapture.

<She is almost as big as Lottie,> someone gasped.

Lottie clapped her hands and danced about, giggling.

<She's dressed for the theatre,> said Lavinia. <Her cloak is lined with ermine.>

<Oh,> cried Ermengarde, darting forward, <she has an opera-glass in her hand—a blue-and-gold one!>

<Here is her trunk,> said Sara. <Let us open it and look at her things.>

She sat down upon the floor and turned the key. The children crowded clamouring around her, as she lifted tray after tray and revealed their contents. Never had the schoolroom been in such an uproar. There were lace collars and silk stockings and handkerchiefs; there was a jewel case containing a necklace and a tiara which looked quite as if they were made of real diamonds; there was a long sealskin and muff, there were ball dresses and walking dresses and visiting dresses; there were hats and tea gowns and fans. Even Lavinia and Jessie forgot that they were too elderly to care for dolls, and uttered exclamations of delight and caught up things to look at them.

<Suppose,> Sara said, as she stood by the table, putting a large, black-velvet hat on the impassively smiling owner of all these splendours—<suppose she understands human talk and feels proud of being admired.>

<You are always supposing things,> said Lavinia, and her air was very superior.

<I know I am,> answered Sara, undisturbedly. <I like it. There is nothing so nice as supposing. It's almost like being a fairy. If you suppose anything hard enough it seems as if it were real.>

<It's all very well to suppose things if you have everything,> said Lavinia. <Could you suppose and pretend if you were a beggar and lived in a garret?>

Sara stopped arranging the Last Doll's ostrich plumes, and looked thoughtful.

<I \textsc{believe} I could,> she said. <If one was a beggar, one would have to suppose and pretend all the time. But it mightn't be easy.>

She often thought afterward how strange it was that just as she had finished saying this—just at that very moment—Miss Amelia came into the room.

<Sara,> she said, <your papa's solicitor, Mr Barrow, has called to see Miss Minchin, and, as she must talk to him alone and the refreshments are laid in her parlour, you had all better come and have your feast now, so that my sister can have her interview here in the schoolroom.>

Refreshments were not likely to be disdained at any hour, and many pairs of eyes gleamed. Miss Amelia arranged the procession into decorum, and then, with Sara at her side heading it, she led it away, leaving the Last Doll sitting upon a chair with the glories of her wardrobe scattered about her; dresses and coats hung upon chair backs, piles of lace-frilled petticoats lying upon their seats.

Becky, who was not expected to partake of refreshments, had the indiscretion to linger a moment to look at these beauties—it really was an indiscretion.

<Go back to your work, Becky,> Miss Amelia had said; but she had stopped to pick up reverently first a muff and then a coat, and while she stood looking at them adoringly, she heard Miss Minchin upon the threshold, and, being smitten with terror at the thought of being accused of taking liberties, she rashly darted under the table, which hid her by its tablecloth.

Miss Minchin came into the room, accompanied by a sharp-featured, dry little gentleman, who looked rather disturbed. Miss Minchin herself also looked rather disturbed, it must be admitted, and she gazed at the dry little gentleman with an irritated and puzzled expression.

She sat down with stiff dignity, and waved him to a chair.

<Pray, be seated, Mr Barrow,> she said.

Mr Barrow did not sit down at once. His attention seemed attracted by the Last Doll and the things which surrounded her. He settled his eyeglasses and looked at them in nervous disapproval. The Last Doll herself did not seem to mind this in the least. She merely sat upright and returned his gaze indifferently.

<A hundred pounds,> Mr Barrow remarked succinctly. <All expensive material, and made at a Parisian modiste's. He spent money lavishly enough, that young man.>

Miss Minchin felt offended. This seemed to be a disparagement of her best patron and was a liberty.

Even solicitors had no right to take liberties.

<I beg your pardon, Mr Barrow,> she said stiffly. <I do not understand.>

<Birthday presents,> said Mr Barrow in the same critical manner, <to a child eleven years old! Mad extravagance, I call it.>

Miss Minchin drew herself up still more rigidly.

<Captain Crewe is a man of fortune,> she said. <The diamond mines alone\longdash>

Mr Barrow wheeled round upon her. <Diamond mines!> he broke out. <There are none! Never were!>

Miss Minchin actually got up from her chair.

<What!> she cried. <What do you mean?>

<At any rate,> answered Mr Barrow, quite snappishly, <it would have been much better if there never had been any.>

<Any diamond mines?> ejaculated Miss Minchin, catching at the back of a chair and feeling as if a splendid dream was fading away from her.

<Diamond mines spell ruin oftener than they spell wealth,> said Mr Barrow. <When a man is in the hands of a very dear friend and is not a businessman himself, he had better steer clear of the dear friend's diamond mines, or gold mines, or any other kind of mines dear friends want his money to put into. The late Captain Crewe\longdash>

Here Miss Minchin stopped him with a gasp.

<The \textsc{late} Captain Crewe!> she cried out. <The \textsc{late}! You don't come to tell me that Captain Crewe is\longdash>

<He's dead, ma'am,> Mr Barrow answered with jerky brusqueness. <Died of jungle fever and business troubles combined. The jungle fever might not have killed him if he had not been driven mad by the business troubles, and the business troubles might not have put an end to him if the jungle fever had not assisted. Captain Crewe is dead!>

Miss Minchin dropped into her chair again. The words he had spoken filled her with alarm.

<What \textsc{were} his business troubles?> she said. <What \textsc{were} they?>

<Diamond mines,> answered Mr Barrow, <and dear friends—and ruin.>

Miss Minchin lost her breath.

<Ruin!> she gasped out.

<Lost every penny. That young man had too much money. The dear friend was mad on the subject of the diamond mine. He put all his own money into it, and all Captain Crewe's. Then the dear friend ran away—Captain Crewe was already stricken with fever when the news came. The shock was too much for him. He died delirious, raving about his little girl—and didn't leave a penny.>

Now Miss Minchin understood, and never had she received such a blow in her life. Her show pupil, her show patron, swept away from the Select Seminary at one blow. She felt as if she had been outraged and robbed, and that Captain Crewe and Sara and Mr Barrow were equally to blame.

<Do you mean to tell me,> she cried out, <that he left \textsc{nothing}! That Sara will have no fortune! That the child is a beggar! That she is left on my hands a little pauper instead of an heiress?>

Mr Barrow was a shrewd businessman, and felt it as well to make his own freedom from responsibility quite clear without any delay.

<She is certainly left a beggar,> he replied. <And she is certainly left on your hands, ma'am—as she hasn't a relation in the world that we know of.>

Miss Minchin started forward. She looked as if she was going to open the door and rush out of the room to stop the festivities going on joyfully and rather noisily that moment over the refreshments.

<It is monstrous!> she said. <She's in my sitting room at this moment, dressed in silk gauze and lace petticoats, giving a party at my expense.>

<She's giving it at your expense, madam, if she's giving it,> said Mr Barrow, calmly. <Barrow \& Skipworth are not responsible for anything. There never was a cleaner sweep made of a man's fortune. Captain Crewe died without paying \textsc{our} last bill—and it was a big one.>

Miss Minchin turned back from the door in increased indignation. This was worse than anyone could have dreamed of its being.

<That is what has happened to me!> she cried. <I was always so sure of his payments that I went to all sorts of ridiculous expenses for the child. I paid the bills for that ridiculous doll and her ridiculous fantastic wardrobe. The child was to have anything she wanted. She has a carriage and a pony and a maid, and I've paid for all of them since the last cheque came.>

Mr Barrow evidently did not intend to remain to listen to the story of Miss Minchin's grievances after he had made the position of his firm clear and related the mere dry facts. He did not feel any particular sympathy for irate keepers of boarding schools.

<You had better not pay for anything more, ma'am,> he remarked, <unless you want to make presents to the young lady. No one will remember you. She hasn't a brass farthing to call her own.>

<But what am I to do?> demanded Miss Minchin, as if she felt it entirely his duty to make the matter right. <What am I to do?>

<There isn't anything to do,> said Mr Barrow, folding up his eyeglasses and slipping them into his pocket. <Captain Crewe is dead. The child is left a pauper. Nobody is responsible for her but you.>

<I am not responsible for her, and I refuse to be made responsible!>

Miss Minchin became quite white with rage.

Mr Barrow turned to go.

<I have nothing to do with that, madam,> he said uninterestedly. <Barrow \& Skipworth are not responsible. Very sorry the thing has happened, of course.>

<If you think she is to be foisted off on me, you are greatly mistaken,> Miss Minchin gasped. <I have been robbed and cheated; I will turn her into the street!>

If she had not been so furious, she would have been too discreet to say quite so much. She saw herself burdened with an extravagantly brought-up child whom she had always resented, and she lost all self-control.

Mr Barrow undisturbedly moved toward the door.

<I wouldn't do that, madam,> he commented; <it wouldn't look well. Unpleasant story to get about in connection with the establishment. Pupil bundled out penniless and without friends.>

He was a clever business man, and he knew what he was saying. He also knew that Miss Minchin was a business woman, and would be shrewd enough to see the truth. She could not afford to do a thing which would make people speak of her as cruel and hard-hearted.

<Better keep her and make use of her,> he added. <She's a clever child, I believe. You can get a good deal out of her as she grows older.>

<I will get a good deal out of her before she grows older!> exclaimed Miss Minchin.

<I am sure you will, ma'am,> said Mr Barrow, with a little sinister smile. <I am sure you will. Good morning!>

He bowed himself out and closed the door, and it must be confessed that Miss Minchin stood for a few moments and glared at it. What he had said was quite true. She knew it. She had absolutely no redress. Her show pupil had melted into nothingness, leaving only a friendless, beggared little girl. Such money as she herself had advanced was lost and could not be regained.

And as she stood there breathless under her sense of injury, there fell upon her ears a burst of gay voices from her own sacred room, which had actually been given up to the feast. She could at least stop this.

But as she started toward the door it was opened by Miss Amelia, who, when she caught sight of the changed, angry face, fell back a step in alarm.

<What \textsc{is} the matter, sister?> she ejaculated.

Miss Minchin's voice was almost fierce when she answered:

<Where is Sara Crewe?>

Miss Amelia was bewildered.

<Sara!> she stammered. <Why, she's with the children in your room, of course.>

<Has she a black frock in her sumptuous wardrobe?>—in bitter irony.

<A black frock?> Miss Amelia stammered again. <A \textsc{black} one?>

<She has frocks of every other colour. Has she a black one?>

Miss Amelia began to turn pale.

<No—ye-es!> she said. <But it is too short for her. She has only the old black velvet, and she has outgrown it.>

<Go and tell her to take off that preposterous pink silk gauze, and put the black one on, whether it is too short or not. She has done with finery!>

Then Miss Amelia began to wring her fat hands and cry.

<Oh, sister!> she sniffed. <Oh, sister! What \textsc{can} have happened?>

Miss Minchin wasted no words.

<Captain Crewe is dead,> she said. <He has died without a penny. That spoiled, pampered, fanciful child is left a pauper on my hands.>

Miss Amelia sat down quite heavily in the nearest chair.

<Hundreds of pounds have I spent on nonsense for her. And I shall never see a penny of it. Put a stop to this ridiculous party of hers. Go and make her change her frock at once.>

<I\@?> panted Miss Amelia. <M-must I go and tell her now?>

<This moment!> was the fierce answer. <Don't sit staring like a goose. Go!>

Poor Miss Amelia was accustomed to being called a goose. She knew, in fact, that she was rather a goose, and that it was left to geese to do a great many disagreeable things. It was a somewhat embarrassing thing to go into the midst of a room full of delighted children, and tell the giver of the feast that she had suddenly been transformed into a little beggar, and must go upstairs and put on an old black frock which was too small for her. But the thing must be done. This was evidently not the time when questions might be asked.

She rubbed her eyes with her handkerchief until they looked quite red. After which she got up and went out of the room, without venturing to say another word. When her older sister looked and spoke as she had done just now, the wisest course to pursue was to obey orders without any comment. Miss Minchin walked across the room. She spoke to herself aloud without knowing that she was doing it. During the last year the story of the diamond mines had suggested all sorts of possibilities to her. Even proprietors of seminaries might make fortunes in stocks, with the aid of owners of mines. And now, instead of looking forward to gains, she was left to look back upon losses.

<The Princess Sara, indeed!> she said. <The child has been pampered as if she were a \textsc{queen}.> She was sweeping angrily past the corner table as she said it, and the next moment she started at the sound of a loud, sobbing sniff which issued from under the cover.

<What is that!> she exclaimed angrily. The loud, sobbing sniff was heard again, and she stooped and raised the hanging folds of the table cover.

<How \textsc{dare} you!> she cried out. <How dare you! Come out immediately!>

It was poor Becky who crawled out, and her cap was knocked on one side, and her face was red with repressed crying.

<If you please, 'm—it's me, mum,> she explained. <I know I hadn't ought to. But I was lookin' at the doll, mum—an' I was frightened when you come in—an' slipped under the table.>

<You have been there all the time, listening,> said Miss Minchin.

<No, mum,> Becky protested, bobbing curtsies. <Not listenin'—I thought I could slip out without your noticin', but I couldn't an' I had to stay. But I didn't listen, mum—I wouldn't for nothin'. But I couldn't help hearin'.>

Suddenly it seemed almost as if she lost all fear of the awful lady before her. She burst into fresh tears.

<Oh, please, 'm,> she said; <I dare say you'll give me warnin', mum—but I'm so sorry for poor Miss Sara—I'm so sorry!>

<Leave the room!> ordered Miss Minchin.

Becky curtsied again, the tears openly streaming down her cheeks.

<Yes, 'm; I will, 'm,> she said, trembling; <but oh, I just wanted to arst you: Miss Sara—she's been such a rich young lady, an' she's been waited on, 'and and foot; an' what will she do now, mum, without no maid? If—if, oh please, would you let me wait on her after I've done my pots an' kettles? I'd do 'em that quick—if you'd let me wait on her now she's poor. Oh,> breaking out afresh, <poor little Miss Sara, mum—that was called a princess.>

Somehow, she made Miss Minchin feel more angry than ever. That the very scullery maid should range herself on the side of this child—whom she realized more fully than ever that she had never liked—was too much. She actually stamped her foot.

<No—certainly not,> she said. <She will wait on herself, and on other people, too. Leave the room this instant, or you'll leave your place.>

Becky threw her apron over her head and fled. She ran out of the room and down the steps into the scullery, and there she sat down among her pots and kettles, and wept as if her heart would break.

<It's exactly like the ones in the stories,> she wailed. <Them pore princess ones that was drove into the world.>

Miss Minchin had never looked quite so still and hard as she did when Sara came to her, a few hours later, in response to a message she had sent her.

Even by that time it seemed to Sara as if the birthday party had either been a dream or a thing which had happened years ago, and had happened in the life of quite another little girl.

Every sign of the festivities had been swept away; the holly had been removed from the schoolroom walls, and the forms and desks put back into their places. Miss Minchin's sitting room looked as it always did—all traces of the feast were gone, and Miss Minchin had resumed her usual dress. The pupils had been ordered to lay aside their party frocks; and this having been done, they had returned to the schoolroom and huddled together in groups, whispering and talking excitedly.

<Tell Sara to come to my room,> Miss Minchin had said to her sister. <And explain to her clearly that I will have no crying or unpleasant scenes.>

<Sister,> replied Miss Amelia, <she is the strangest child I ever saw. She has actually made no fuss at all. You remember she made none when Captain Crewe went back to India. When I told her what had happened, she just stood quite still and looked at me without making a sound. Her eyes seemed to get bigger and bigger, and she went quite pale. When I had finished, she still stood staring for a few seconds, and then her chin began to shake, and she turned round and ran out of the room and upstairs. Several of the other children began to cry, but she did not seem to hear them or to be alive to anything but just what I was saying. It made me feel quite queer not to be answered; and when you tell anything sudden and strange, you expect people will say \textsc{something}—whatever it is.>

Nobody but Sara herself ever knew what had happened in her room after she had run upstairs and locked her door. In fact, she herself scarcely remembered anything but that she walked up and down, saying over and over again to herself in a voice which did not seem her own, <My papa is dead! My papa is dead!>

Once she stopped before Emily, who sat watching her from her chair, and cried out wildly, <Emily! Do you hear? Do you hear—papa is dead? He is dead in India—thousands of miles away.>

When she came into Miss Minchin's sitting room in answer to her summons, her face was white and her eyes had dark rings around them. Her mouth was set as if she did not wish it to reveal what she had suffered and was suffering. She did not look in the least like the rose-coloured butterfly child who had flown about from one of her treasures to the other in the decorated schoolroom. She looked instead a strange, desolate, almost grotesque little figure.

She had put on, without Mariette's help, the cast-aside black-velvet frock. It was too short and tight, and her slender legs looked long and thin, showing themselves from beneath the brief skirt. As she had not found a piece of black ribbon, her short, thick, black hair tumbled loosely about her face and contrasted strongly with its pallor. She held Emily tightly in one arm, and Emily was swathed in a piece of black material.

<Put down your doll,> said Miss Minchin. <What do you mean by bringing her here?>

<No,> Sara answered. <I will not put her down. She is all I have. My papa gave her to me.>

She had always made Miss Minchin feel secretly uncomfortable, and she did so now. She did not speak with rudeness so much as with a cold steadiness with which Miss Minchin felt it difficult to cope—perhaps because she knew she was doing a heartless and inhuman thing.

<You will have no time for dolls in future,> she said. <You will have to work and improve yourself and make yourself useful.>

Sara kept her big, strange eyes fixed on her, and said not a word.

<Everything will be very different now,> Miss Minchin went on. <I suppose Miss Amelia has explained matters to you.>

<Yes,> answered Sara. <My papa is dead. He left me no money. I am quite poor.>

<You are a beggar,> said Miss Minchin, her temper rising at the recollection of what all this meant. <It appears that you have no relations and no home, and no one to take care of you.>

For a moment the thin, pale little face twitched, but Sara again said nothing.

<What are you staring at?> demanded Miss Minchin, sharply. <Are you so stupid that you cannot understand? I tell you that you are quite alone in the world, and have no one to do anything for you, unless I choose to keep you here out of charity.>

<I understand,> answered Sara, in a low tone; and there was a sound as if she had gulped down something which rose in her throat. <I understand.>

<That doll,> cried Miss Minchin, pointing to the splendid birthday gift seated near—<that ridiculous doll, with all her nonsensical, extravagant things—I actually paid the bill for her!>

Sara turned her head toward the chair.

<The Last Doll,> she said. <The Last Doll.> And her little mournful voice had an odd sound.

<The Last Doll, indeed!> said Miss Minchin. <And she is mine, not yours. Everything you own is mine.>

<Please take it away from me, then,> said Sara. <I do not want it.>

If she had cried and sobbed and seemed frightened, Miss Minchin might almost have had more patience with her. She was a woman who liked to domineer and feel her power, and as she looked at Sara's pale little steadfast face and heard her proud little voice, she quite felt as if her might was being set at naught.

<Don't put on grand airs,> she said. <The time for that sort of thing is past. You are not a princess any longer. Your carriage and your pony will be sent away—your maid will be dismissed. You will wear your oldest and plainest clothes—your extravagant ones are no longer suited to your station. You are like Becky—you must work for your living.>

To her surprise, a faint gleam of light came into the child's eyes—a shade of relief.

<Can I work?> she said. <If I can work it will not matter so much. What can I do?>

<You can do anything you are told,> was the answer. <You are a sharp child, and pick up things readily. If you make yourself useful I may let you stay here. You speak French well, and you can help with the younger children.>

<May I\@?> exclaimed Sara. <Oh, please let me! I know I can teach them. I like them, and they like me.>

<Don't talk nonsense about people liking you,> said Miss Minchin. <You will have to do more than teach the little ones. You will run errands and help in the kitchen as well as in the schoolroom. If you don't please me, you will be sent away. Remember that. Now go.>

Sara stood still just a moment, looking at her. In her young soul, she was thinking deep and strange things. Then she turned to leave the room.

<Stop!> said Miss Minchin. <Don't you intend to thank me?>

Sara paused, and all the deep, strange thoughts surged up in her breast.

<What for?> she said.

<For my kindness to you,> replied Miss Minchin. <For my kindness in giving you a home.>

Sara made two or three steps toward her. Her thin little chest heaved up and down, and she spoke in a strange un-childishly fierce way.

<You are not kind,> she said. <You are \textsc{not} kind, and it is \textsc{not} a home.> And she had turned and run out of the room before Miss Minchin could stop her or do anything but stare after her with stony anger.

She went up the stairs slowly, but panting for breath and she held Emily tightly against her side.

<I wish she could talk,> she said to herself. <If she could speak—if she could speak!>

She meant to go to her room and lie down on the tiger-skin, with her cheek upon the great cat's head, and look into the fire and think and think and think. But just before she reached the landing Miss Amelia came out of the door and closed it behind her, and stood before it, looking nervous and awkward. The truth was that she felt secretly ashamed of the thing she had been ordered to do.

<You—you are not to go in there,> she said.

<Not go in?> exclaimed Sara, and she fell back a pace.

<That is not your room now,> Miss Amelia answered, reddening a little.

Somehow, all at once, Sara understood. She realized that this was the beginning of the change Miss Minchin had spoken of.

<Where is my room?> she asked, hoping very much that her voice did not shake.

<You are to sleep in the attic next to Becky.>

Sara knew where it was. Becky had told her about it. She turned, and mounted up two flights of stairs. The last one was narrow, and covered with shabby strips of old carpet. She felt as if she were walking away and leaving far behind her the world in which that other child, who no longer seemed herself, had lived. This child, in her short, tight old frock, climbing the stairs to the attic, was quite a different creature.

When she reached the attic door and opened it, her heart gave a dreary little thump. Then she shut the door and stood against it and looked about her.

Yes, this was another world. The room had a slanting roof and was whitewashed. The whitewash was dingy and had fallen off in places. There was a rusty grate, an old iron bedstead, and a hard bed covered with a faded coverlet. Some pieces of furniture too much worn to be used downstairs had been sent up. Under the skylight in the roof, which showed nothing but an oblong piece of dull gray sky, there stood an old battered red footstool. Sara went to it and sat down. She seldom cried. She did not cry now. She laid Emily across her knees and put her face down upon her and her arms around her, and sat there, her little black head resting on the black draperies, not saying one word, not making one sound.

And as she sat in this silence there came a low tap at the door—such a low, humble one that she did not at first hear it, and, indeed, was not roused until the door was timidly pushed open and a poor tear-smeared face appeared peeping round it. It was Becky's face, and Becky had been crying furtively for hours and rubbing her eyes with her kitchen apron until she looked strange indeed.

<Oh, miss,> she said under her breath. <Might I—would you allow me—jest to come in?>

Sara lifted her head and looked at her. She tried to begin a smile, and somehow she could not. Suddenly—and it was all through the loving mournfulness of Becky's streaming eyes—her face looked more like a child's not so much too old for her years. She held out her hand and gave a little sob.

<Oh, Becky,> she said. <I told you we were just the same—only two little girls—just two little girls. You see how true it is. There's no difference now. I'm not a princess anymore.>

Becky ran to her and caught her hand, and hugged it to her breast, kneeling beside her and sobbing with love and pain.

<Yes, miss, you are,> she cried, and her words were all broken. <Whats'ever 'appens to you—whats'ever—you'd be a princess all the same—an' nothin' couldn't make you nothin' different.>